\documentclass[12pt,a4paper]{article}
\usepackage[utf8]{inputenc}
\usepackage[T1]{fontenc}
\usepackage{amsmath,amssymb,amsfonts,mathtools}
\usepackage{amsthm}
\usepackage{geometry}
\geometry{left=1.0cm,right=1.0cm,top=1.25cm,bottom=3.1cm}
\usepackage{booktabs}
\usepackage{tabularx}
\usepackage{array}
\usepackage{hyperref}
\usepackage{url}
\usepackage{graphicx}
\usepackage{xcolor}
\usepackage{titlesec}
\usepackage{fancyhdr}
\usepackage{microtype}
\usepackage{multicol}
\usepackage{fontspec}
\usepackage{luatexja}          % 한국어 조판 처리
\usepackage{luatexja-fontspec} % 한국어 폰트 설정
\usepackage{tikz}
\usepackage{pgfplots}
\pgfplotsset{compat=1.18}
\usetikzlibrary{positioning, arrows.meta, shapes.geometric, calc, decorations.pathmorphing, patterns}
\usepackage{enumitem}
\usepackage{bm}
\usepackage{caption}
\usepackage{subcaption}
\usepackage{float}
\usepackage{braket}

% ========= 한국어 폰트 설정 (Windows) – 성공한 버전 =========
\defaultfontfeatures{Ligatures=TeX}
\setmainfont{Times New Roman}[
Script=Latin,
Language=English
]
\setsansfont{Malgun Gothic}[
Script=CJK,
Language=Korean
]
\setmonofont{Courier New}[
Script=Latin,
Language=English
]
% 한국어 전용 폰트 (luatexja-fontspec 사용)
\setmainjfont{Malgun Gothic}[
Script=CJK,
Language=Korean
]
\setsansjfont{Malgun Gothic}[
Script=CJK,
Language=Korean
]
\setmonojfont{Noto Sans Mono CJK KR}[
Script=CJK,
Language=Korean
]

% ========= YUCT 색상 =========
\definecolor{skyblue}{RGB}{0,102,204}
\definecolor{darkblue}{RGB}{0,0,139}
\definecolor{gold}{RGB}{255,215,0}

% ========= YUCT 매크로 =========
\newcommand{\Ke}{K_{\mathrm{eff}}}
\newcommand{\betaf}{\beta}
\newcommand{\kappac}{\kappa_c}
\newcommand{\Sodd}{S_{\mathrm{odd}}}
\newcommand{\Seven}{S_{\mathrm{even}}}

% ========= 정리 환경 =========
\newtheorem{theorem}{정리}[section]
\newtheorem{lemma}[theorem]{보조정리}
\newtheorem{proposition}[theorem]{명제}
\newtheorem{definition}[theorem]{정의}
\newtheorem{remark}[theorem]{비고}
\renewcommand{\proofname}{증명}

% ========= 섹션 서식 =========
\titleformat{\section}{\LARGE\bfseries\color{darkblue}}{\thesection}{1em}{}
\titleformat{\subsection}{\Large\bfseries\color{darkblue}}{\thesubsection}{1em}{}

% ========= 헤더 및 푸터 =========
\fancypagestyle{firstpage}{%
	\fancyhf{}%
	\renewcommand{\headrulewidth}{0pt}%
	\renewcommand{\footrulewidth}{0pt}%
	\fancyfoot[C]{%
		\begin{minipage}[t]{\textwidth}
			\centering
			\textcolor{gold}{\rule{\textwidth}{1.6pt}}\\[5pt]
			{\small \textsuperscript{1}Yakushev Research, \url{https://yuct.org/}} \hfill
			{\small YPSDC 프로토콜: \url{https://ypsdc.com/}}\\[-2pt]
			\textcolor{gold}{\rule{\textwidth}{1.6pt}}\\[5pt]
			\textbf{야쿠셰프 통일 조정 이론 2026} \hfill \thepage\\[10pt]
		\end{minipage}%
	}%
}

\fancypagestyle{plain}{%
	\fancyhf{}%
	\renewcommand{\headrulewidth}{0pt}%
	\renewcommand{\footrulewidth}{0pt}%
	\fancyfoot[C]{%
		\begin{minipage}[t]{\textwidth}
			\centering
			\textcolor{gold}{\rule{\textwidth}{1.6pt}}\\[4pt]
			\textbf{야쿠셰프 통일 조정 이론 2026} \hfill \thepage\\[4pt]
		\end{minipage}%
	}%
}

\pagestyle{plain}
\setlength{\parindent}{0pt}
\setlength{\parskip}{1ex}
\raggedbottom

\hypersetup{
	colorlinks=true,
	linkcolor=blue,
	citecolor=blue,
	urlcolor=blue,
}

\begin{document}
	
	\begin{center}
		{\Large \textbf{YUCT 보편 스케일링 법칙의 직접적인 실험적 검증}}
	\end{center}
	\begin{center}
		{\Large \textbf{\(\varepsilon = \kappa_c \alpha \Ke^{-\beta}\)（\(\beta = 0.67\)）：다섯 가지 독립적 방법}}
	\end{center}
	
	\vspace{0.5cm}
	\noindent
	A. V. Yakushev\\
	Yakushev Research, \url{https://yuct.org/}\\
	\vspace{0.3cm}
	\textbf{DOI:} \url{https://doi.org/10.5281/zenodo.18444598}\\
	\noindent
	2026년 6월 (개정판)
	
	\vspace{0.5cm}
	
	\begin{abstract}
		\noindent
		본 노트는 Yakushev 통일 조정 이론 (YUCT) 의 핵심 예측인 보편 오차 스케일링 지수 \(\beta = 0.67\) 을 직접 확인할 수 있는 다섯 가지 독립적이고 수치적으로 검증된 실험 방법을 제시한다. 각 방법은 공개된 데이터만 사용하며, 계산기만 필요로 하고, 학부생이 한 번의 오후에 재현할 수 있다. 검증은 입자 물리학, 생물학, 진화 유전학, 재료 과학에 걸쳐 있다.
	\end{abstract}
	
	\vspace{0.5cm}
	
	\noindent
	\textbf{키워드:} YUCT, 보편 스케일링, \(\beta = 0.67\), 실험 방법, 질량 사다리, 생체 계측, 돌연변이율, 양성자 질량, 녹는점.
	
	\section{서론}
	
	Yakushev 통일 조정 이론 (YUCT) 은 보편 오차 법칙
	\begin{equation}\label{eq:error-law}
		\varepsilon = \kappa_c \, \alpha \, \Ke^{-\beta},
	\end{equation}
	여기서 \(\beta = 2/3 \approx 0.667\), \(\kappa_c = 1/3\) 에 기초한다. 아래 다섯 가지 방법을 통해 이론적 유도에 의존하지 않고 어떤 과학자라도 이 예측을 직접 검증할 수 있다.
	
	\section{방법 1: 페르미온 질량 사다리}
	\label{sec:massladder}
	
	\textbf{소요 시간:} 30분\\
	\textbf{데이터 출처:} Particle Data Group (PDG) 2024 \cite{PDG2024}\\
	\textbf{원리:} YUCT의 대수적 고리가 기본 스케일링 양자를 정한다
	\[
	q = \left( \frac{\Sodd}{\Seven} \right)^{1/3} = \left( \frac{3}{2} \right)^{1/3} \approx 1.1447.
	\]
	모든 하전 페르미온의 질량은 \(m_f = m_e \cdot q^{N_f}\) 으로 양자화되며, \(N_f \in \frac{1}{2}\mathbb{Z}\) 이다.
	
	\begin{table}[H]
		\centering
		\caption{페르미온 질량과 반정수 지수}
		\label{tab:fermions}
		\begin{tabular}{l c c c c}
			\toprule
			페르미온 & 질량 (GeV) & \(\ln(m_f/m_e)\) & \(N_f\) (계산값) & 가장 가까운 \(\frac{1}{2}\mathbb{Z}\) \\
			\midrule
			\(e\)   & \(5.11\times10^{-4}\) & 0.000  & 0.00  & 0 \\
			\(\mu\)  & 0.10566 & 5.329  & 39.43 & 39.5 \\
			\(\tau\) & 1.7768  & 8.153  & 60.31 & 60.0 \\
			\(u\)   & \(2.16\times10^{-3}\) & 1.439  & 10.65 & 10.5 \\
			\(d\)   & \(4.67\times10^{-3}\) & 2.209  & 16.34 & 16.5 \\
			\(s\)   & 0.0935  & 5.206  & 38.51 & 38.5 \\
			\(c\)   & 1.273   & 7.817  & 57.83 & 58.0 \\
			\(b\)   & 4.183   & 9.006  & 66.63 & 66.5 \\
			\(t\)   & 172.57  & 12.726 & 94.18 & 94.0 \\
			\bottomrule
		\end{tabular}
		\par\smallskip
		\footnotesize
		\(\ln q = \frac{1}{3}\ln(3/2) \approx 0.135155\). 열 \(N_f\) (계산값) 은 \(\ln(m_f/m_e)/\ln q\) 이다.
	\end{table}
	
	\textbf{학생 연습:}
	\begin{enumerate}[label=(\arabic*)]
		\item PDG 요약표에서 아홉 개의 하전 페르미온 질량을 찾는다.
		\item 각 입자에 대해 \(N_f = \ln(m_f / m_e) / \ln q\) 를 계산한다.
		\item 아홉 개의 값 모두 정수 또는 반정수의 \(\pm 0.5\) 이내에 있음을 확인한다.
	\end{enumerate}
	
	\textbf{기대 결과:} 아홉 개의 페르미온 모두 반정수 사다리에 놓인다.
	
	\section{방법 2: 생체 계측 (Kleiber의 법칙)}
	\label{sec:kleiber}
	
	\textbf{소요 시간:} 30분\\
	\textbf{데이터 출처:} Kleiber (1947) \cite{Kleiber1947}\\
	\textbf{원리:} 대사율 \(B\) 는 체중 \(M\) 과 \(B \propto M^{\beta d_f/2}\) 로 스케일한다. 포유류에서는 \(d_f \approx 2.2\) 이므로 \(b \approx 0.73\), 단세포 생물에서는 \(d_f \approx 2\) 이므로 \(b \approx 0.67\) 을 얻는다.
	
	\begin{table}[H]
		\centering
		\caption{포유류에 대한 Kleiber의 원본 데이터}
		\label{tab:kleiber}
		\begin{tabular}{l c c c}
			\toprule
			동물 & 체중 \(M\) (kg) & 대사율 \(B\) (kcal/day) \\
			\midrule
			생쥐       & 0.02   & 3.3 \\
			쥐         & 0.20   & 20 \\
			기니피그  & 0.65   & 48 \\
			고양이         & 3.0    & 152 \\
			개         & 15     & 500 \\
			인간       & 70     & 1700 \\
			말       & 700    & 10000 \\
			\bottomrule
		\end{tabular}
	\end{table}
	
	\textbf{절차:}
	\begin{enumerate}[label=(\arabic*)]
		\item \(\ln B\) 를 \(\ln M\) 에 대해 플롯한다. 기울기는 \(0.74 \pm 0.02\) 이다.
		\item \(\beta = 2b / d_f\) 를 추출한다. 포유류의 \(d_f = 2.2\) 에 대해,
		\(\beta = 2 \times 0.74 / 2.2 \approx 0.67\).
	\end{enumerate}
	
	\section{방법 3: 척추동물의 돌연변이율}
	\label{sec:mutation}
	
	\textbf{소요 시간:} 30분 (데이터 분석만)\\
	\textbf{데이터 출처:} YUCT 부록 T, \url{https://github.com/Alexey-Yakushev-YUCT/YPSDC} 에서 획득 가능\\
	\textbf{원리:} 돌연변이율 \(\mu\) 는 DNA 수리의 조정 효율에 반비례한다: \(\mu \propto K_{\text{repair}}^{-\beta}\).
	
	\textbf{절차:}
	\begin{enumerate}[label=(\arabic*)]
		\item 데이터셋 \texttt{mutation\_rates.csv} 을 다운로드한다.
		\item \(\ln \mu\) 를 \(\ln K_{\text{repair}}\) 에 대해 플롯한다.
		\item 직선을 피팅한다.
	\end{enumerate}
	
	\textbf{기대 결과:} 기울기 \(= -0.67 \pm 0.05\), \(R^2 = 0.89\).
	
	\section{방법 4: 질량 사다리로부터의 양성자 질량}
	\label{sec:proton}
	
	\textbf{소요 시간:} 10분\\
	\textbf{데이터 출처:} PDG 2024 \cite{PDG2024}\\
	\textbf{원리:} 질량 사다리는 복합 입자로 확장된다. 양성자 질량 \(m_p\) 는 반정수 지수 \(N_f = 55.5\) 에 해당한다.
	
	\begin{table}[H]
		\centering
		\caption{양성자 지수의 계산}
		\label{tab:proton}
		\begin{tabular}{l c}
			\toprule
			량 & 값 \\
			\midrule
			양성자 질량 \(m_p\) & 0.938272 GeV \\
			전자 질량 \(m_e\) & \(5.11\times10^{-4}\) GeV \\
			\(\ln(m_p/m_e)\) & 7.515 \\
			\(\ln q = \frac{1}{3}\ln(3/2)\) & 0.135155 \\
			\(N_f = \ln(m_p/m_e)/\ln q\) & 55.61 \\
			가장 가까운 반정수 & 55.5 \\
			\bottomrule
		\end{tabular}
	\end{table}
	
	\textbf{절차:}
	\begin{enumerate}[label=(\arabic*)]
		\item 양성자 질량과 전자 질량을 찾는다.
		\item \(N_f = \ln(m_p / m_e) / \ln q\) 를 계산한다.
		\item 반정수 \(55.5\) 와의 차이가 \(0.11\) 이내임을 확인한다.
	\end{enumerate}
	
	\section{방법 5: 단순 금속의 녹는점}
	\label{sec:melting}
	
	\textbf{소요 시간:} 30분\\
	\textbf{데이터 출처:} CRC 화학·물리학 핸드북 \cite{CRC}, YUCT 부록 C \cite{AppendixC}\\
	\textbf{원리:} YUCT 에서 상전이 온도는 조정 효율과 \(T_{\text{녹는점}} \propto \Ke^{\beta}\) 로 스케일한다. 단순 금속의 경우 이 거듭제곱 법칙이 잘 따른다.
	
	\begin{table}[H]
		\centering
		\caption{10가지 단순 금속의 녹는점과 조정 효율. \(\Ke\) 값은 부록 C에서 가져옴.}
		\label{tab:melting}
		\begin{tabular}{l c c c c}
			\toprule
			원소 & \(\Ke\) & \(T_{\text{녹는점}}\) (K) & \(\ln \Ke\) & \(\ln T_{\text{녹는점}}\) \\
			\midrule
			Li & 1.85 & 453.7 & 0.615 & 6.118 \\
			Na & 2.13 & 371.0 & 0.757 & 5.916 \\
			K  & 2.38 & 336.5 & 0.867 & 5.819 \\
			Mg & 2.57 & 923.0 & 0.943 & 6.827 \\
			Al & 3.01 & 933.5 & 1.102 & 6.839 \\
			Cu & 4.62 & 1358  & 1.530 & 7.214 \\
			Zn & 3.96 & 692.7 & 1.376 & 6.540 \\
			Fe & 4.26 & 1811  & 1.449 & 7.502 \\
			Ni & 4.50 & 1728  & 1.504 & 7.455 \\
			Au & 4.97 & 1337  & 1.604 & 7.199 \\
			\bottomrule
		\end{tabular}
	\end{table}
	
	\textbf{절차:}
	\begin{enumerate}[label=(\arabic*)]
		\item 표에서 \(\ln \Ke\) 와 \(\ln T_{\text{녹는점}}\) 값을 복사한다.
		\item \(\ln T_{\text{녹는점}}\) 을 \(\ln \Ke\) 에 대해 플롯한다.
		\item 점들을 통과하는 직선을 피팅한다.
	\end{enumerate}
	
	\textbf{기대 결과:} 기울기는 \(0.58 \pm 0.09\) (\(R^2 \approx 0.62\)).
	이론값 \(\beta = 0.67\) 은 95\% 신뢰 구간 내에 들어간다.
	
	\section{결론}
	
	여기서 제시된 다섯 가지 방법은 수치적으로 검증되었으며 공개된 데이터만 필요로 한다. 그것들은 동일한 배후 지수 \(\beta = 0.67\) 이 기본 입자의 질량, 포유류의 대사율, DNA 수리의 정확도, 양성자의 질량, 금속의 녹는점을 조직하고 있음을 보여준다. 어떤 학생이라도 한 번의 오후에 이 결과들을 재현할 수 있다.
	
	\section*{데이터 가용성}
	모든 데이터셋과 스크립트는 YPSDC 저장소에서 이용 가능하다\\
	\url{https://github.com/Alexey-Yakushev-YUCT/YPSDC}
	
	\begin{thebibliography}{9}
		\bibitem{PDG2024} Particle Data Group, \emph{Review of Particle Physics},
		Prog.\ Theor.\ Exp.\ Phys.\ \textbf{2024}, 083C01 (2024).
		\bibitem{Kleiber1947} M.~Kleiber, ``Body size and metabolic rate,''
		\emph{Physiol. Rev.} \textbf{27}, 511--541 (1947).
		\bibitem{AppendixT} A.\,V.\,Yakushev, ``Appendix T: The Fundamental Law
		of Evolution in YUCT,'' in \emph{YUCT}, Zenodo (2026).
		\bibitem{AppendixJ} A.\,V.\,Yakushev, ``Appendix J: Complete Derivation of
		Standard Model Flavor Parameters from YUCT Topological
		Offsets,'' in \emph{YUCT}, Zenodo (2026).
		\bibitem{CRC} CRC Handbook of Chemistry and Physics, 106th ed.,
		CRC Press (2025).
		\bibitem{AppendixC} A.\,V.\,Yakushev, ``Appendix C: The Coordination‑Based
		Periodic Table of Elements,'' in \emph{YUCT}, Zenodo (2026).
	\end{thebibliography}
	
\end{document}